% GENERATED FILE. Edit canonical lessons, then run npm run generate:books.
% canonical-source-hash: fnv1a64:79a55bdff524ee68
% canonical-lessons: GE-C389-leisten, GE-C389-lohnen, GE-C389-finanzieren, GE-C389-ausgabe, GE-C389-einnahme

\chapter{Sich etwas leisten}
\label{ch:ge-sich-etwas-leisten}
\hlchaptermodality{389}

\begin{chapteropening}
\textbf{By the end of this chapter:} \emph{I can say what I can afford, whether something is worth it, how I pay for it, and what goes out and comes in.}

Complete the last lesson of chapter 389: I can say what I can afford, whether something is worth it, how I pay for it, and what goes out and comes in.
\end{chapteropening}

\section[leisten]{leisten — to afford (sich leisten); to achieve}

\label{lesson:GE-C389-leisten}



\begin{quote}
Before the new one: say the German for to pay back, then the German for to borrow.
\end{quote}

\subsection*{You'll want to know: leisten}
\textbf{leisten} — \textquotedblleft{}to afford (sich leisten); to achieve\textquotedblright{}.

\begin{quote}
Ein neues Auto können wir uns im Moment nicht leisten.
\end{quote}

\begin{quote}
\emph{We cannot afford a new car at the moment.}
\end{quote}

\emph{Sich etwas leisten} — to afford something. Also to achieve: \emph{Sie hat viel geleistet.} \emph{Die Leistung} — the achievement.

\begin{cousinweb}
\emph{leisten} is inherited from Old High German \emph{leisten} and Proto-Germanic \emph{laistijaną} (EtymDB), usually explained as first \textquotedblleft{}to follow a track\textquotedblright{}, then to carry out what you owe; EtymDB links it to English \emph{last}.
\end{cousinweb}

\begin{grammarlens}[title={What you've built}]
One more word for telling, explaining and describing at B1.
\end{grammarlens}

\subsection*{Your turn}
\begin{itemize}
\raggedright
  \item \emph{Say it:} \emph{leisten}
  \item \emph{Say it:} \emph{leisten}, once more
  \item \emph{Say it:} a sentence of your own with \emph{leisten}
  \item \emph{Recall:} say the German for to pay back, then the German for to borrow, then say \emph{leisten} again
\end{itemize}

\begin{tcolorbox}[breakable,colback=teal!4,colframe=teal!35!black,title={Before you move on}]
What does leisten mean? (\textquotedblleft{}To afford (sich leisten); to achieve\textquotedblright{}.) Say it once more.
\end{tcolorbox}
\section[lohnen]{lohnen — to be worth it (sich lohnen); to reward}

\label{lesson:GE-C389-lohnen}



\begin{quote}
Before the new one: say the German for to borrow, then the German for to afford.
\end{quote}

\subsection*{You'll want to know: lohnen}
\textbf{lohnen} — \textquotedblleft{}to be worth it (sich lohnen); to reward\textquotedblright{}.

\begin{quote}
Der lange Weg hat sich gelohnt: Die Aussicht ist fantastisch.
\end{quote}

\begin{quote}
\emph{The long walk was worth it: the view is fantastic.}
\end{quote}

Mostly reflexive: \emph{Das lohnt sich} — it is worth it. \emph{Der Lohn} — the wage, the reward.

\begin{cousinweb}
\emph{lohnen} is inherited from Middle High German \emph{lōnen} and Old High German \emph{lōnōn} \textquotedblleft{}to pay, to reward\textquotedblright{} (EtymDB), from \emph{Lohn} \textquotedblleft{}reward\textquotedblright{}, Proto-Germanic \emph{launą}.
\end{cousinweb}

\begin{grammarlens}[title={What you've built}]
One more word for telling, explaining and describing at B1.
\end{grammarlens}

\subsection*{Your turn}
\begin{itemize}
\raggedright
  \item \emph{Say it:} \emph{lohnen}
  \item \emph{Say it:} \emph{lohnen}, once more
  \item \emph{Say it:} a sentence of your own with \emph{lohnen}
  \item \emph{Recall:} say the German for to borrow, then the German for to afford, then say \emph{lohnen} again
\end{itemize}

\begin{tcolorbox}[breakable,colback=teal!4,colframe=teal!35!black,title={Before you move on}]
What does lohnen mean? (\textquotedblleft{}To be worth it (sich lohnen); to reward\textquotedblright{}.) Say it once more.
\end{tcolorbox}
\section[finanzieren]{finanzieren — to finance, to pay for}

\label{lesson:GE-C389-finanzieren}



\begin{quote}
Before the new one: say the German for to afford, then the German for to be worth it.
\end{quote}

\subsection*{You'll want to know: finanzieren}
\textbf{finanzieren} — \textquotedblleft{}to finance, to pay for\textquotedblright{}.

\begin{quote}
Sie finanziert ihr Studium mit einem Nebenjob.
\end{quote}

\begin{quote}
\emph{She pays for her studies with a part-time job.}
\end{quote}

Perfekt: \emph{hat finanziert}. \emph{Die Finanzen} — finances.

\begin{cousinweb}
\emph{finanzieren} was borrowed from French \emph{financer} (EtymDB), from \emph{finance} \textquotedblleft{}payment\textquotedblright{}, built on Old French \emph{finer} \textquotedblleft{}to settle a debt\textquotedblright{}, from \emph{fin} \textquotedblleft{}end\textquotedblright{}, Latin \emph{fīnis} (Lewis \& Short): you bring a debt to an end.
\end{cousinweb}

\begin{grammarlens}[title={What you've built}]
One more word for telling, explaining and describing at B1.
\end{grammarlens}

\subsection*{Your turn}
\begin{itemize}
\raggedright
  \item \emph{Say it:} \emph{finanzieren}
  \item \emph{Say it:} \emph{finanzieren}, once more
  \item \emph{Say it:} a sentence of your own with \emph{finanzieren}
  \item \emph{Recall:} say the German for to afford, then the German for to be worth it, then say \emph{finanzieren} again
\end{itemize}

\begin{tcolorbox}[breakable,colback=teal!4,colframe=teal!35!black,title={Before you move on}]
What does finanzieren mean? (\textquotedblleft{}To finance, to pay for\textquotedblright{}.) Say it once more.
\end{tcolorbox}
\section[die Ausgabe]{die Ausgabe — the expense; the edition}

\label{lesson:GE-C389-ausgabe}



\begin{quote}
Before the new one: say the German for to be worth it, then the German for to finance.
\end{quote}

\subsection*{You'll want to know: die Ausgabe}
\textbf{die Ausgabe} — \textquotedblleft{}the expense; the edition\textquotedblright{}.

\begin{quote}
Unsere größte Ausgabe ist die Miete.
\end{quote}

\begin{quote}
\emph{Our biggest expense is the rent.}
\end{quote}

Plural: \emph{die Ausgaben}. Also an edition of a newspaper or a book: \emph{die neue Ausgabe}.

\begin{cousinweb}
\emph{Ausgabe} is inherited from Middle High German \emph{ūzgābe} (EtymDB), of the family of \emph{ausgeben} \textquotedblleft{}to spend, to give out\textquotedblright{}, built on \emph{geben} \textquotedblleft{}to give\textquotedblright{}, Proto-Germanic \emph{gebaną}.
\end{cousinweb}

\begin{grammarlens}[title={What you've built}]
One more word for telling, explaining and describing at B1.
\end{grammarlens}

\subsection*{Your turn}
\begin{itemize}
\raggedright
  \item \emph{Say it:} \emph{die Ausgabe}
  \item \emph{Say it:} \emph{die Ausgabe}, once more
  \item \emph{Say it:} a sentence of your own with \emph{die Ausgabe}
  \item \emph{Recall:} say the German for to be worth it, then the German for to finance, then say \emph{die Ausgabe} again
\end{itemize}

\begin{tcolorbox}[breakable,colback=teal!4,colframe=teal!35!black,title={Before you move on}]
What does die Ausgabe mean? (\textquotedblleft{}The expense; the edition\textquotedblright{}.) Say it once more.
\end{tcolorbox}
\section[die Einnahme]{die Einnahme — the income, the takings}

\label{lesson:GE-C389-einnahme}



\begin{quote}
Before the new one: say the German for to finance, then the German for the expense.
\end{quote}

\subsection*{You'll want to know: die Einnahme}
\textbf{die Einnahme} — \textquotedblleft{}the income, the takings\textquotedblright{}.

\begin{quote}
Die Einnahmen des Cafés sind im Sommer viel höher.
\end{quote}

\begin{quote}
\emph{The café's takings are much higher in summer.}
\end{quote}

Mostly plural: \emph{die Einnahmen}. Also the taking of medicine: \emph{die Einnahme von Tabletten}.

\begin{cousinweb}
\emph{Einnahme} is the noun of \emph{einnehmen} \textquotedblleft{}to take in\textquotedblright{}, built on \emph{nehmen} \textquotedblleft{}to take\textquotedblright{}, Old High German \emph{neman}, Proto-Germanic \emph{nemaną} (EtymDB).
\end{cousinweb}

\begin{grammarlens}[title={What you've built}]
One more word for telling, explaining and describing at B1.
\end{grammarlens}

\subsection*{Your turn}
\begin{itemize}
\raggedright
  \item \emph{Say it:} \emph{die Einnahme}
  \item \emph{Say it:} \emph{die Einnahme}, once more
  \item \emph{Say it:} a sentence of your own with \emph{die Einnahme}
  \item \emph{Recall:} say the German for to finance, then the German for the expense, then say \emph{die Einnahme} again
\end{itemize}

\begin{tcolorbox}[breakable,colback=teal!4,colframe=teal!35!black,title={Before you move on}]
What does die Einnahme mean? (\textquotedblleft{}The income, the takings\textquotedblright{}.) Say it once more.
\end{tcolorbox}
